\chapter{Le principe : trois dossiers, un contrat}

Ce chapitre est la carte. Tout le reste en découle, et si vous ne deviez lire
qu'un chapitre, ce serait celui-ci.

\section{Trois dossiers, trois natures de contribution}

Vous ne manipulez \textbf{jamais} de DLL et ne lancez \textbf{jamais} de commande
de build. Vous déposez un fichier dans un dossier, et l'atelier fait le reste.

\begin{center}
\begin{tabular}{@{}p{3.4cm}p{4.6cm}p{5.8cm}@{}}
\toprule
\textbf{J'ajoute\ldots} & \textbf{Je dépose\ldots} & \textbf{Ce qui se passe} \\
\midrule
un moteur de règles & \texttt{rules/mes\_regles.cpp} & détecté → \textbf{compilé} → chargé → panneau \emph{Modules} \\
une IA & \texttt{ai/mon\_ia.cpp} & idem → liste des pilotes, panneau \emph{Joueurs} \\
une grille & \texttt{boards/diamant.json} & \textbf{aucune compilation} → panneau \emph{Règles} \\
\bottomrule
\end{tabular}
\end{center}

La différence entre les deux premiers et le troisième n'est pas une commodité,
c'est une décision de conception.

\begin{nkretenir}[Une règle est du code, un plateau est une donnée]
\begin{quote}
\emph{Le plateau n'est pas une constante du code : c'est un descripteur
sérialisable, éditable à la souris dans le moteur de test.}
\hfill\textsc{regles\_completes\_v2.md:95}
\end{quote}
Une règle décrit un \emph{comportement} : elle se programme. Un plateau décrit
une \emph{forme} : il se décrit. Confondre les deux vous condamne à recompiler
pour essayer un hexagone au lieu d'un carré — et donc à ne jamais l'essayer.

Le chapitre 5 montre la troisième voie : un plateau décrit \textbf{en C++}, sans
JSON, quand sa forme ne se laisse pas mettre en tableau.
\end{nkretenir}

\section{Le cycle de travail}

\begin{nkcode}{Ce qui se passe quand vous sauvegardez}
1.  vous ecrivez   travail/rules/mes_regles.cpp
2.  vous sauvegardez
3.  l'atelier detecte le changement (il scrute une fois par seconde)
4.  il COMPILE  ->  il charge  ->  votre module apparait dans le menu
5.  vous rejouez une partie, sans avoir quitte l'application
\end{nkcode}

En cas d'erreur, \textbf{la sortie complète du compilateur} s'affiche dans le
panneau \emph{Modules}. C'est le seul retour que vous aurez : pas de terminal,
pas de chaîne de build à apprendre.

\begin{nkpiege}[La compilation fige l'atelier une à deux secondes]
C'est assumé, et documenté dans \nkfile{src/ConquerorLab/NkcSession.h} :

\begin{quote}\small
La compilation est SYNCHRONE : quand le stagiaire sauvegarde, l'atelier se fige
une seconde ou deux, puis son module apparaît. C'est assumé — un compilateur
piloté depuis un thread annexe demanderait une file de travaux et un verrou sur
le catalogue, pour gagner deux secondes toutes les cinq minutes.
\end{quote}

Si l'atelier se fige au moment où vous sauvegardez, il ne plante pas : il vous
compile.
\end{nkpiege}

\section{Le contrat : deux tables de pointeurs}

Un module ne dérive d'aucune classe et n'exporte aucun symbole C++. Il remplit
une \textbf{structure de pointeurs de fonction} et l'expose par deux fonctions
\texttt{C}.

\begin{nkcode}{include/Conqueror/ConquerorRulesABI.h:255-273 (extrait)}
struct NkcRulesFactory {
        NkcRulesInfo   info;
        NkcRulesVTable vtable;
};

// Les DEUX symboles que tout module de regles doit exporter.
using NkcRulesGetFactoryFn   = void (*)(NkcRulesFactory *out);
using NkcRulesSetAllocatorFn = void (*)(NkcAllocFn a, NkcFreeFn f);

#define NKC_RULES_SYM_GET_FACTORY "nkc_rules_get_factory"
#define NKC_RULES_SYM_SET_ALLOC   "nkc_rules_set_allocator"
\end{nkcode}

Pourquoi une table de pointeurs plutôt qu'une classe abstraite ? Parce qu'une
classe abstraite impose une \textbf{ABI C++} — disposition de vtable, décoration
des noms, modèle d'exception — qui change d'un compilateur à l'autre et parfois
d'une version à l'autre. Une structure de pointeurs de fonction, non.

\section{Qui parle à qui}

\begin{nkcode}{Le schema qui explique le decoupage des deux stages}
                    +---------------------------+
                    |        ConquerorLab       |
                    | (l'atelier - six panneaux)|
                    +-------------+-------------+
                                  | NE CONNAIT AUCUNE REGLE
                    +-------------+-------------+
                    |                           |
          NkcRulesVTable                 NkcAIVTable
                    |                           |
       +------------v-----------+   +-----------v------------+
       | votre moteur de regles |<--|        votre IA        |
       |          (A1)          |   |          (A2)          |
       +------------------------+   +------------------------+
                  ^                              |
                  +------------------------------+
                     l'IA TRAVERSE les regles :
                     elle recoit la vtable + un handle opaque
\end{nkcode}

Trois propriétés découlent de ce schéma, et ce sont elles qui font tenir le
stage :

\begin{enumerate}
  \item \textbf{A2 démarre en semaine 1}, sur un module bouchon, sans attendre A1 ;
  \item \textbf{quand A1 change son code interne, l'IA continue de fonctionner} —
        elle n'a jamais vu ce code ;
  \item \textbf{les deux ne peuvent pas diverger sur les règles} : il n'en existe
        qu'une implémentation, et l'IA la traverse.
\end{enumerate}

\section{L'état est opaque, la vue est en lecture seule}

Vous choisissez librement la représentation interne de votre partie. Ni l'atelier
ni l'IA ne la voient. Ce qu'ils voient est une \textbf{vue}.

\begin{nkcode}{include/Conqueror/ConquerorRulesABI.h:92-109 (abrege)}
struct NkcStateView {
        const NkcCellView *cells     = nullptr;  // cellCount entrees
        const NkcCoord    *coords    = nullptr;  // meme indexation
        uint32             cellCount = 0;

        NkcTopology topology    = NkcTopology::HexPointy;
        uint8       playerCount = 2;
        uint8       current     = 0;   // joueur au trait
        uint8       finished    = 0;
        int8        winner      = -2;  // -1 = nul, -2 = en cours
        uint32      turn        = 0;

        const int32 *energy         = nullptr;
        const int32 *conquestTenths = nullptr;   // en DIXIEMES
        const int32 *totemCount     = nullptr;
};
\end{nkcode}

\begin{nkpiege}[Les pointeurs de la vue appartiennent au module]
Ils restent valides \textbf{jusqu'au prochain appel mutant} sur cet état
(\texttt{ApplyMove}, \texttt{Setup}, \texttt{DeserializeState},
\texttt{CloneState}). L'appelant ne libère jamais rien.

En pratique : après avoir appliqué un coup, \textbf{redemandez la vue}. Les
champs scalaires (\texttt{current}, \texttt{turn}, \texttt{finished}) sont des
\emph{copies} faites au moment de \texttt{GetView} — ils ne se mettent pas à jour
tout seuls.
\end{nkpiege}

\section{Les quatre exigences non négociables}

Elles ne sont pas des préférences de style. Elles conditionnent l'existence même
du sujet A2 et de toute campagne de mesure.

\begin{nkcode}{REGLES\_COMPLETES\_v2.md:583-590}
1. Zero flottant dans la logique de regles. Entiers partout, PC en dixiemes.
2. PRNG porte par l'etat de partie, serialise avec lui. Jamais de generateur
   global.
3. Rejeu bit-a-bit : seed + liste de coups -> etat final identique, sur
   Windows et sur Linux, quel que soit le compilateur.
4. Aucun ordre d'iteration de conteneur ne doit etre observable dans le
   resultat. Toute resolution ambigue est arbitree par un critere explicite
   (index de joueur, coordonnee), jamais par l'ordre de parcours d'une table.
\end{nkcode}

Chacune a une raison très concrète.

\subsection*{1 — Zéro flottant}

Les Points de Conquête valent $1 + niveau \times 0{,}5 + ennemis \times 0{,}2 +
alli\acute{e}s \times 0{,}1$. En virgule flottante, \textbf{l'ordre
d'accumulation change le dernier bit} selon le compilateur et la plateforme, et
le rejeu bit-à-bit tombe. Le moteur stocke donc les PC en \textbf{dixièmes
entiers} : base 10, niveau $\times 5$, ennemi $\times 2$, allié $\times 1$. La
division par 10 est un fait d'\emph{affichage} — elle n'existe nulle part dans la
logique.

\subsection*{2 — PRNG porté par l'état}

L'IA clone une position des milliers de fois par seconde. Avec un générateur
global, deux clones piochent dans le même flux et le rejeu meurt. Au palier 0 il
n'y a aucun aléatoire dans Conqueror ; le champ existe quand même, parce que
l'ajouter plus tard obligerait à re-sérialiser tous les états déjà enregistrés.

\subsection*{3 — Rejeu bit-à-bit}

C'est ce qui permet de coller un journal dans un rapport de bug et de le rejouer
à l'identique ailleurs. Le panneau \emph{Journal} affiche l'empreinte après
\textbf{chaque} coup : quand deux rejeux divergent, elle dit \emph{où}, pas
seulement \emph{qu'ils diffèrent}.

\subsection*{4 — Aucun ordre observable}

Si votre générateur de coups parcourt une table de hachage, l'ordre des coups
change entre deux exécutions. L'IA choisit alors un coup différent, la partie
diverge, et vos dix mille parties ne mesurent plus rien.

\begin{nkretenir}
\texttt{GenerateLegalMoves} \textbf{doit} produire les coups dans un ordre stable
et reproductible. Dans les exemples de ce cours : cases par index croissant, puis
voisins dans l'ordre déclaré. Ne le changez jamais sans raison — et si vous le
changez, sachez que toutes les mesures antérieures deviennent incomparables.
\end{nkretenir}

\section{Ce qu'un module a le droit d'inclure}

La frontière est \textbf{NKCanvas / NKGui} : tout ce qui est en dessous vous est
ouvert, rien de ce qui est au-dessus ne l'est.

\begin{nkcode}{src/ConquerorLab/NkcModuleCompiler.h — StackIncludes}
/// L'atelier ne donnait au depart que trois -I et aucun lien : un
/// module compilait sans rien lier du moteur. C'etait elegant, et
/// trop etroit -- un moteur de regles a besoin d'une chaine, d'un
/// tableau dynamique, d'un formatage, d'un journal. Reecrire tout
/// cela a la main dans chaque module est du temps vole au jeu.
\end{nkcode}

Vous avez donc droit à :

\begin{center}
\begin{tabular}{@{}p{7.4cm}p{6.4cm}@{}}
\toprule
\textbf{Ce que vous incluez} & \textbf{Ce que ça vous donne} \\
\midrule
\nkfile{Conqueror/ConquerorRulesABI.h} & le contrat \\
\nkfile{Conqueror/ConquerorGeometry.h} & voisinage, distance — inline et entiers \\
\nkfile{NKCore/NkTypes.h} & \texttt{uint32}, \texttt{int8}, \texttt{usize} \\
\nkfile{NKContainers/String/NkString.h} & \texttt{NkString}, \texttt{NkFormat} \\
\nkfile{NKContainers/Sequential/NkVector.h} & \texttt{NkVector}, \texttt{NkList} \\
\nkfile{NKContainers/Associative/NkMap.h} & \texttt{NkMap}, \texttt{NkSet} \\
\nkfile{NKMath/NkFunctions.h} & \texttt{math::NkSqrt}, \texttt{NkClamp}\ldots \\
\nkfile{NKLogger/NkLog.h} & \texttt{logger.Infof(...)} \\
\nkfile{NKFileSystem/NkFile.h} & \texttt{NkFile}, \texttt{NkDirectory} \\
\nkfile{NKThreading/NkThread.h} & \texttt{NkThread}, \texttt{NkMutex} \\
la bibliothèque standard C & \texttt{<cstring>}, \texttt{<cstdio>}\ldots \\
\bottomrule
\end{tabular}
\end{center}

Vous n'avez \textbf{pas} droit à NKCanvas, NKGui, NKWindow, NKRenderer — un
moteur de règles ne dessine pas. C'est l'atelier qui affiche, à partir de la vue
que vous exposez. \textbf{Cette frontière est le cœur du contrat}, pas une
restriction administrative : c'est elle qui permet de changer entièrement
l'affichage sans toucher une règle, et inversement.

\section{Quand le contrat change : majeure et mineure}

Le contrat va évoluer pendant vos huit semaines. La question qui compte est :
\textbf{faut-il tout recompiler à chaque fois ?}

\begin{nkcode}{include/Conqueror/ConquerorRulesABI.h}
inline constexpr uint32 kRulesAbiMajor = 3;
inline constexpr uint32 kRulesAbiMinor = 1;
\end{nkcode}

\begin{itemize}
  \item \textbf{MAJEURE} — un champ change de type ou de place, une signature
        change. Votre module est \textbf{refusé} : il faut recompiler, et le
        panneau \emph{Modules} le dit.
  \item \textbf{MINEURE} — une fonction est ajoutée \emph{à la fin} de la vtable.
        Votre module \textbf{continue de tourner}. L'atelier l'accepte et note :
        « module en ABI 3.0, atelier en 3.1 — il fonctionne, les fonctions
        ajoutées depuis ne sont pas disponibles. »
\end{itemize}

Pourquoi cette distinction existe : sans elle, le jour où j'ai ajouté trois
entrées à la vtable, \textbf{tous} les modules devenaient « ABI 2, attendue 3 » —
refusés alors qu'ils tournaient parfaitement. Vous faire recompiler tout votre
travail parce que \emph{j'ai} ajouté une fonction est exactement la friction
qu'un contrat doit supprimer.

\begin{nkpiege}[Ce que ça a coûté d'apprendre]
La règle « on ajoute à la fin » est plus étroite qu'elle n'en a l'air : on ajoute
à la fin de la \textbf{dernière structure de la chaîne}.

\texttt{NkcRulesInfo} précède \texttt{NkcRulesVTable} dans la fabrique. Grossir
\texttt{info} \textbf{décale} la vtable : un module d'époque écrit sa table à
l'ancien décalage, l'atelier la lit au nouveau, et on appelle des adresses au
hasard. Premier essai : \textbf{segfault}. Les champs sont donc à la fin de
\texttt{NkcRulesFactory}, après la vtable.

\nkfile{tests/NkcAbiCompat.cpp} monte désormais la garde : il compile un module
contre des en-têtes \textbf{figés} et le fait jouer avec l'atelier d'aujourd'hui.
\textbf{11 vérifications, 0 échec.}
\end{nkpiege}

\section{Afficher un état pour comprendre — le panneau Sortie}

Votre module a \textbf{sa propre copie de la pile} : l'édition de liens est
statique. Conséquence directe, et surprenante la première fois : un
\texttt{logger.Infof()} depuis votre code n'atteint pas la console de l'atelier.
Vous écrivez dans \emph{votre} journal, pas dans le sien.

Ce n'est pas un oubli — c'est la conséquence de l'isolation qui fait tout
l'intérêt du système. Mais déboguer un moteur de règles sans pouvoir afficher un
état, c'est déboguer à l'aveugle. L'atelier \textbf{injecte donc un puits} dans
chaque module au chargement, et le branchement tient en une ligne :

\begin{nkcode}{exemples/rules/RegleContratNu.cpp — fin de fichier}
NKC_MODULE_EXPORT void nkc_rules_set_allocator(NkcAllocFn a, NkcFreeFn f) {
    gAlloc = a; gFree = f;
}
NKC_MODULE_LOGGING(rules)   // <- une ligne, et NKC_LOG_* va dans « Sortie »
\end{nkcode}

\texttt{NKC\_MODULE\_LOGGING(rules)} dans un module de règles,
\texttt{NKC\_MODULE\_LOGGING(ai)} dans une IA. Il faut inclure
\nkfile{Conqueror/ConquerorLog.h}. À partir de là :

\begin{nkcode}{exemples/rules/RegleContratNu.cpp — V\_Setup}
NKC_LOG_INFO("nouvelle partie : %d cases, %d joueurs, graine %llu",
             kCells, (int32)s->playerCount, (unsigned long long)seed);
\end{nkcode}

s'affiche dans le panneau \textbf{Sortie}, avec le nom de votre module et le
niveau.

\begin{nkretenir}[\texttt{logger.Infof} marche aussi]
Si vous incluez NKLogger, la macro branche en plus un \texttt{NkISink} sur le
logger \emph{privé} de votre module, qui repousse tout vers l'atelier. Le réflexe
de quelqu'un qui connaît déjà Nkentseu fonctionne donc sans qu'il ait à
l'apprendre.

Une seule contrainte : incluez \nkfile{ConquerorLog.h} \textbf{après} NKLogger —
la détection se fait sur son garde d'inclusion.
\end{nkretenir}

\begin{nkpiege}[Deux limites à connaître]
\textbf{Sans atelier, tout part sur stderr.} Banc d'essai, test en ligne de
commande : personne n'injecte de puits, et vos traces sortent sous la forme
\texttt{[INFO] module: ...}. Elles ne disparaissent jamais en silence.

\textbf{Ne journalisez pas dans une boucle chaude.} \texttt{NKC\_LOG\_*} formate
\emph{avant} de savoir si quelqu'un écoute. Dans un rollout MCTS c'est un
désastre — et le tampon de l'atelier est borné à 4096 lignes, donc vous perdriez
de toute façon ce que vous cherchez. Une trace par coup, pas par nœud. Pour
mesurer, il y a \texttt{NkcAIResult} et \texttt{GetDebugJson}.
\end{nkpiege}

L'isolation explique aussi \nkfile{nkc_rules_set_allocator} : sans injection, les
deux côtés allouent dans des tas distincts.

Les trois exemples de ce cours n'utilisent volontairement que
\texttt{<cstring>} et \texttt{<cstdio>} : ils doivent rester lisibles par
quelqu'un qui ne connaît pas encore Nkentseu. Rien ne vous oblige à cette
austérité.

\begin{nkpiege}[Pourquoi la projection écran n'est pas dans la géométrie]
\begin{quote}\small
Aucune dépendance à NKMath : tout y est inline et entier, donc un module compilé
par l'atelier n'a AUCUN symbole à lier. La projection écran (cos/sin/arrondi) vit
ailleurs — c'est de la présentation, pas de la règle.
\hfill\textsc{ConquerorGeometry.h:10-13}
\end{quote}

\texttt{NkRound} de NKMath est un symbole \emph{lié}. L'inclure dans le contrat
ferait échouer l'édition de liens de \textbf{tous} les modules.
\end{nkpiege}

\section{La géométrie fournie, en trois fonctions}

C'est le socle commun. Le chapitre 5 montre comment s'en passer entièrement.

\begin{nkcode}{include/Conqueror/ConquerorGeometry.h:26-88 (signatures)}
int32    NeighborCount(NkcTopology t) noexcept;
NkcCoord Neighbor(NkcTopology t, NkcCoord c, int32 i) noexcept;
int32    Distance(NkcTopology t, NkcCoord a, NkcCoord b) noexcept;
\end{nkcode}

Quatre topologies, et le voisinage est une propriété de la topologie — jamais
codé en dur ailleurs :

\begin{center}
\begin{tabular}{@{}lll@{}}
\toprule
\textbf{Topologie} & \textbf{Voisins} & \textbf{Coordonnées} \\
\midrule
\texttt{HexPointy} & 6 & axiales $(q, r)$, pointes haut/bas \\
\texttt{HexFlat}   & 6 & axiales $(q, r)$, pointes gauche/droite \\
\texttt{Square4}   & 4 & $(x, y)$, orthogonaux \\
\texttt{Square8}   & 8 & $(x, y)$, $+$ diagonales \\
\bottomrule
\end{tabular}
\end{center}

Un seul type de coordonnée pour les quatre : c'est la topologie qui décide de
l'interprétation.

\begin{nkpiege}[L'ordre des voisins est NORMATIF]
\begin{quote}\small
i-ème voisin de \texttt{c}. L'ORDRE EST NORMATIF : il fixe l'ordre de génération
des coups, donc la reproductibilité. Ne jamais le changer.
\hfill\textsc{ConquerorGeometry.h:36-37}
\end{quote}
\end{nkpiege}

\begin{nkexo}[1 — Lire le contrat]
Ouvrez \nkfile{ConquerorRulesABI.h} et faites la liste des fonctions de
\texttt{NkcRulesVTable} \textbf{sans regarder le chapitre 2}. Classez-les en
quatre familles : cycle de vie, réglages, jeu, sérialisation. Combien y en a-t-il ?
Laquelle vous paraît la plus difficile à implémenter, et pourquoi ?
\end{nkexo}

\begin{nkexo}[2 — La cascade]
Lisez \nkfile{REGLES\_COMPLETES\_v2.md} \S7.3. Une duplication retourne les
ennemis voisins de la case où l'on vient de poser. Si un totem retourné a
lui-même des voisins ennemis, sont-ils retournés à leur tour ? Trouvez la phrase
qui tranche, et expliquez pourquoi le document dit que c'est « le point le plus
facile à implémenter de travers ».
\end{nkexo}

\begin{nkexo}[3 — Le flottant qui tue]
Écrivez un petit programme qui additionne \texttt{0.1f} cent fois, puis fait la
même chose en dixièmes entiers, et compare au résultat exact \texttt{10.0}. Quel
écart obtenez-vous ? Multipliez-le mentalement par dix mille parties et deux
plateformes : vous tenez la raison de l'exigence n\textsuperscript{o} 1.
\end{nkexo}
